%────────────────────── Capítulo 2 — Descripción ───────────────────────
% Detalles sobre el mecanismo de ataque y su implementación.
\chapter{Descripción del ataque}
\label{cap:descripcion}

\section{AES-128 y el valor intermedio atacado}
\label{sec:aes}

AES-128 cifra bloques de 16 bytes con una clave de 16 bytes
\cite{nist2023fips197}. El bloque se copia en un \emph{estado} de 16 bytes, al
que se aplica un primer \texttt{AddRoundKey} con la clave y, después, diez
rondas. Cada ronda encadena cuatro transformaciones: \texttt{SubBytes}, que
sustituye cada byte por su imagen en una tabla fija de 256 entradas (la S-box);
\texttt{ShiftRows}, que rota las filas del estado; \texttt{MixColumns}, que
combina linealmente los bytes de cada columna, y \texttt{AddRoundKey}, que hace
el XOR del estado con una subclave. La última ronda omite \texttt{MixColumns}.
Las once subclaves se obtienen de la clave mediante la expansión de clave. La
S-box es pública y no tiene nada secreto; es el único paso no lineal del
cifrado.

Para el ataque solo importan los dos primeros pasos, que trabajan byte a byte.
El \texttt{AddRoundKey} inicial usa la clave tal cual, así que el byte $j$ del
estado pasa a valer $p_j \oplus k_j$, donde $p_j$ es el byte $j$ del texto
claro y $k_j$ el de la clave. El \texttt{SubBytes} de la primera ronda lo
transforma en
\begin{equation}
  v_j = \mathrm{SBox}(p_j \oplus k_j),
  \label{eq:valor-intermedio}
\end{equation}
que es el valor intermedio que se ataca en este trabajo. Hay dos razones para
elegirlo.

La primera es que depende de un solo byte de clave y de un byte que el atacante
conoce, porque es quien elige el texto claro. Para cada uno de los 256 valores
posibles de $k_j$ se puede calcular qué valor $v_j$ habría procesado el
dispositivo, y comprobar cuál de esas predicciones encaja con el consumo medido.
Cada byte se ataca por separado: en lugar de $2^{128}$ claves hay que
comprobar $16 \times 256 = 4096$ hipótesis.

La segunda explica por qué se ataca la salida de la S-box y no la del XOR, que
también depende solo de $p_j$ y $k_j$. Con el modelo de fuga del peso de
Hamming (\cref{sec:cpa}), la predicción del candidato complementario
$k \oplus \mathtt{0xff}$ sería $\operatorname{HW}(p_j \oplus k \oplus
\mathtt{0xff}) = 8 - \operatorname{HW}(p_j \oplus k)$: exactamente la contraria,
con la misma correlación en valor absoluto que el candidato correcto, y no
habría forma de distinguirlos. Los candidatos que difieren del correcto en un
solo bit también predecirían casi lo mismo. La S-box rompe esa estructura,
porque dos entradas parecidas dan salidas que no se parecen, y solo el candidato
correcto predice bien la fuga.

El mismo razonamiento vale para la última ronda, que no tiene
\texttt{MixColumns}: el byte del criptograma $c_j$ es la salida de la S-box,
desplazada por \texttt{ShiftRows}, combinada por XOR con un byte de la subclave
de la ronda 10. La hipótesis $\mathrm{InvSBox}(c_j \oplus k)$ permite atacar
esa subclave, y como la expansión de clave de AES-128 es invertible, de ella se
deduce la clave. En este trabajo el ataque principal es el de la primera ronda;
el de la última se usa en la simulación como comprobación adicional
(\cref{sec:validacion-simulada}).

\section{Análisis simple de consumo (SPA)}
\label{sec:spa}

El análisis simple consiste en observar una traza, o la media de unas pocas, y
reconocer en ella las operaciones del algoritmo \cite[sec.~II.B.1]{mishra2025modern}.
Cuando el camino de ejecución depende de la clave, como en una exponenciación
modular de RSA que solo multiplica cuando un bit vale uno, el SPA puede revelar
la clave directamente. En el artículo que presentó el DPA ya se distinguían a
simple vista las 16 rondas de un DES en una sola traza \cite{kocher1999dpa}.

Un AES software bien escrito ejecuta las mismas instrucciones sea cual sea la
clave, así que el SPA no basta para atacarlo. En este trabajo sirve para otra
cosa: localizar la primera ronda dentro de la traza, para capturar solo esa
ventana y no buscar a ciegas, y comprobar que el montaje funciona. Que se
distingan las diez rondas en una traza es el primer criterio de validación de la
captura, y todavía no se ha comprobado.
% TODO: traza real con las rondas marcadas cuando el montaje funcione (fase 4).

\section{Análisis diferencial de consumo (DPA)}
\label{sec:dpa}

El DPA de Kocher, Jaffe y Jun \cite{kocher1999dpa} fue el primer ataque
estadístico sobre el consumo. Se miden $N$ cifrados con textos claros aleatorios
y se guarda, para cada uno, la traza $T_i$ y el texto claro $p_i$. Para cada
byte $j$ y cada candidato $k$, una \emph{función de selección} predice un bit
del valor intermedio; aquí, el bit $\ell$ de $\mathrm{SBox}(p_{i,j} \oplus k)$.
Las trazas se reparten en dos grupos, $G_0$ y $G_1$, según el valor predicho de
ese bit, y se restan sus medias en cada instante:
\begin{equation}
  \Delta_j(k, t) = \frac{1}{|G_1|} \sum_{i \in G_1} T_i[t]
                 - \frac{1}{|G_0|} \sum_{i \in G_0} T_i[t].
  \label{eq:dpa}
\end{equation}

Si el candidato es falso, la partición no tiene relación con lo que hace el
dispositivo, los dos grupos son estadísticamente iguales y la diferencia tiende
a cero al crecer $N$. Si es el correcto, en el instante en que se procesa $v_j$
uno de los grupos consume sistemáticamente algo más que el otro y aparece un
pico. El candidato elegido es el que da el pico más alto.

El DPA de un bit desaprovecha información: el dispositivo filtra los ocho bits
de $v_j$ a la vez, y los siete que no se usan para partir se comportan como
ruido. Además, el bit más informativo depende de la implementación y el
atacante no sabe de antemano cuál escoger.

\section{Análisis de consumo por correlación (CPA)}
\label{sec:cpa}

El CPA de Brier, Clavier y Olivier \cite{brier2004cpa} sustituye la partición
por un \emph{modelo de fuga} que predice el consumo entero, y la diferencia de
medias por la correlación de Pearson
\cite[sec.~II.B.3]{mishra2025modern}.

El modelo más habitual es el peso de Hamming $\operatorname{HW}(v)$, el número
de bits a uno del valor procesado, que supone que el consumo en el instante en
que se manipula $v$ es aproximadamente $a \operatorname{HW}(v) + b$ más ruido.
Brier et al.\ proponen en general la distancia de Hamming entre el valor y un
estado de referencia ---el contenido previo del registro o del bus---, del que
el peso de Hamming es el caso con referencia nula. Aquí se usa el peso de
Hamming porque no se conoce qué valor ocupaba antes el registro en el
microcontrolador. Es una hipótesis de trabajo: si las medidas encajan mejor con
otro modelo, se cambiará.

Para cada byte $j$, cada candidato $k$ y cada instante $t$ se calcula la
correlación entre la fuga predicha $h_i = \operatorname{HW}(\mathrm{SBox}(p_{i,j}
\oplus k))$ y la muestra medida $T_i[t]$ sobre las $N$ trazas:
\begin{equation}
  \rho_j(k, t) =
  \frac{\sum_{i=1}^{N} (h_i - \bar h)\,(T_i[t] - \bar T[t])}
       {\sqrt{\sum_{i=1}^{N} (h_i - \bar h)^2}\,\sqrt{\sum_{i=1}^{N} (T_i[t] - \bar T[t])^2}},
  \label{eq:cpa}
\end{equation}
donde $\bar h$ y $\bar T[t]$ son las medias sobre las trazas. El candidato
elegido para el byte $j$ es
\begin{equation}
  \hat k_j = \arg\max_{k} \, \max_{t} \, \lvert \rho_j(k, t) \rvert .
  \label{eq:cpa-decision}
\end{equation}
Se usa el valor absoluto porque el signo de la fuga depende del montaje; por
ejemplo, de en qué lado de la alimentación se coloca el \emph{shunt}.

Con el candidato correcto, la predicción coincide con la fuga real y la
correlación sube en el instante en que se calcula su S-box. Con los otros 255,
la predicción es casi independiente de la traza y la correlación se queda en el
nivel del azar, que baja como $1/\sqrt{N}$. La correlación del candidato
correcto, en cambio, no crece con $N$: depende solo de la relación entre la
varianza de la fuga y la del ruido. La clave aparece cuando el máximo de los
falsos cae por debajo del correcto, y por eso el número de trazas necesario
crece aproximadamente como el inverso del cuadrado de la correlación del
correcto. La \cref{sec:validacion-simulada} cuantifica esta relación sobre
trazas simuladas.

Frente al DPA de un bit, el CPA aprovecha los ocho bits del valor intermedio y
da una medida normalizada, comparable entre bytes y entre montajes. El coste de
cálculo es moderado: $16 \times 256$ correlaciones por instante de la traza,
que se pueden calcular de una vez como un producto de matrices. Además, como la
predicción solo depende del byte de texto claro, basta acumular para cada byte
$j$ y cada uno de los 256 valores de $p_j$ cuántas trazas hay y su suma, sin
guardar las trazas; así lo hace la web explicativa del proyecto.

\section{Implementación experimental}
\label{sec:implementacion}

Esta sección describe el montaje previsto. El análisis ya está implementado y
validado en simulación, pero el instrumental está pendiente de confirmar y el
firmware y la captura todavía no existen. Lo que se dice aquí sobre el
hardware es un diseño, no un resultado: se confirmará o se corregirá con las
medidas del \cref{cap:resultados}.

\subsection{Arquitectura}

El montaje sigue el esquema habitual de los laboratorios de análisis de consumo
sobre firmware, como los de ChipWhisperer \cite{cw-lab-dpa,cw-lab-cpa}, pero con
componentes de propósito general:

\begin{enumerate}
  \item El PC envía por serie un texto claro aleatorio de 16 bytes al
    microcontrolador.
  \item El firmware activa una salida digital de disparo, cifra con la clave de
    laboratorio, desactiva el disparo y devuelve el criptograma.
  \item Un osciloscopio mide en un canal la caída de tensión en una resistencia
    de \emph{shunt} en serie con la alimentación del microcontrolador, que es
    proporcional a su corriente, y en otro la señal de disparo.
  \item El PC comprueba el criptograma con su propia implementación de AES y
    guarda la terna (texto claro, criptograma, traza).
\end{enumerate}

Se repite el proceso miles de veces y las trazas se atacan con el mismo código
que se ha validado en simulación.
% TODO: esquema del montaje en figuras/ cuando esté decidido.

\subsection{Microcontrolador y firmware}

El dispositivo atacado es un ATmega328P, el microcontrolador de 8 bits del
Arduino Nano. Está por decidir si se usa la placa o el chip suelto en una placa
de prototipos. En la placa, el regulador, el conversor USB-serie y los
condensadores de desacoplo comparten alimentación con el microcontrolador, y es
difícil medir solo su corriente; el chip suelto, con su cristal y sus
condensadores, permite colocar el \emph{shunt} solo en su línea de alimentación
o de masa, con el mismo firmware.
% TODO: decisión Nano / chip suelto (AGENTS.md, hardware/README.md).

El firmware implementará un AES-128 software sin protecciones, con
\texttt{SubBytes} como una consulta a tabla byte a byte. Ese recorrido secuencial
del estado es el que hace que cada byte del valor intermedio se procese en su
propia ventana temporal, que es la hipótesis del modelo de trazas simuladas. Se
verificará con los ejemplos del FIPS-197 \cite{nist2023fips197} antes de medir.
% TODO: fase 3. Protocolo serie, velocidad, pin de disparo y commit del
% firmware cuando existan.

\subsection{Medida del consumo}

La corriente del microcontrolador se convierte en tensión con una resistencia
de \emph{shunt} de bajo valor. Si se coloca en la línea de masa, la caída se
puede medir respecto a masa con una sonda normal, aunque hay que vigilar los
bucles de masa entre el PC, la placa y el osciloscopio. Los condensadores de
desacoplo filtran precisamente las variaciones rápidas de corriente que se
quieren medir, así que su posición respecto al \emph{shunt} es una decisión del
diseño.
% TODO: valor del shunt, posición y necesidad de amplificador (preguntas 3 y 4
% de docs/reunion-2026-10-19.md).

El osciloscopio puede ser uno del laboratorio de la universidad o uno USB
comprado; está pendiente de que lo confirme un profesor. Tiene que cumplir tres
requisitos: muestrear lo bastante rápido para resolver la actividad de cada
instrucción, disparar con la señal del firmware o capturarla en un segundo canal
para alinear las trazas después, y poder controlarse desde Python en Linux para
automatizar la campaña. Si el reloj del microcontrolador es demasiado rápido
para el muestreo disponible, se puede reducir, a costa de reajustar la velocidad
del puerto serie.
% TODO: modelo de osciloscopio, frecuencia de muestreo y reloj del MCU, con
% enlace y fecha de consulta de las especificaciones.

\subsection{Procesamiento de las trazas}

El análisis está escrito en Python, con \texttt{numpy} para el cálculo y
\texttt{matplotlib} para las figuras, en el directorio \texttt{analisis/} del
repositorio. Incluye una implementación propia de AES-128, verificada con los
vectores del FIPS-197, de la que salen las predicciones del valor intermedio; el
CPA de la \cref{eq:cpa} y el DPA de la \cref{eq:dpa}, vectorizados sobre los 256
candidatos y todas las muestras; las métricas que se reportan (correlación
frente al tiempo, pico por candidato, puesto de la clave correcta frente al
número de trazas) y el generador de trazas sintéticas de la
\cref{sec:validacion-simulada}. Con una semilla fija, la salida es idéntica
entre ejecuciones, y la integración continua lo comprueba en cada cambio.

Las trazas reales se guardarán en un fichero \texttt{.npz} por campaña con las
trazas, los textos claros, los criptogramas y la clave, más un \texttt{.json}
de metadatos con la frecuencia de muestreo, el reloj del microcontrolador, el
\emph{shunt}, el osciloscopio, la fecha y la versión del firmware. Antes del
ataque habrá que alinear las trazas con la señal de disparo y, si es necesario,
recortarlas a la ventana de la primera ronda que se identifique por SPA.
% TODO: alineación y preprocesado reales cuando haya trazas (fase 4).
